\documentclass{ltjbook}
\usepackage{docmute} 
\usepackage{../../myStyle} 

\begin{document}
ここまで尤度を使った推定方法について説明してきました。モーメント法と最尤法，いずれも同じ答えが出るのですがその考え方が違うというところを確認してきましたね。さてここからは，もう1つの推定法であるベイズ法について解説します。そのためには，ベイズ法の根本的な法則について知っておいた方が良いので，少し回りくどいかもしれませんが，まずはそこから説明していきましょう。目的を先に教えてくれないとやだよ，という人のために一言で目的だけ言っておきます。ベイズ法は最尤法と同じく確率モデルを使うのですが，尤度が「確率じゃない」という妙な性質があったのに対し，ベイズ法は答えを確率で出します。尤度を確率にするための仕掛けがこれから始まるのです。

\section{ベイズの定理}
ベイズというのは人の名前で，トーマス・ベイズ（Thomas Bayes、1702年 - 1761年4月17日）はイギリスの長老派の牧師であり数学者でもある人です。この人の考えた「逆確率の法則」から端を発した推定法が，21世紀の現代においてとても重要なものになってきています。ベイズ牧師の考えついた法則は非常に単純なもので，簡単に書くと次のような式で表現できます。
\[ P(A|B) = \frac{P(B|A)P(A)}{P(B)} \]

この式の意味するところを，確率の用語をいろいろ紹介しながら説明していきたいと思います。

話をわかりやすくするために，具体的な数値例とともに話を進めましょう。
たとえば，ある学校の生徒数は男子が40人，女子が60人だったとします。そこからランダムに一人を選び出した時，それが男子である確率はどれぐらいでしょうか？ この確率を$P(\mbox{男子})$とすると，
\[P(\mbox{男子})=\frac{40}{(100)}=0.4 \]
であることはすぐにわかりますね。

では次に\keyterm{同時確率}{どうじかくりつ}{joint probability}という考え方を導入しましょう。これは先ほどの学校について，表\ref{tbl::24_02}のように学年情報を追加した例を使って説明します。

\begin{table}[ht]
    \centering
    \caption{ある学校の学年と性別のデータ} 
    \begin{tabular}{rrrrr}
      \hline
      & 1年 & 2年 & 3年 & 合計 \\
    \hline
    男子 & 15 & 12 & 13 & 40\\
    女子 & 22 & 20 & 18 & 60\\\hline
    合計 & 37 & 32 & 31 & 100\\\hline
    \end{tabular}
      \label{tbl::24_02}
\end{table}

記号で表現するために，性別の変数をAとし，$A_1=$男子，$A_2=$女子，学年の変数をBとして$B_1=$1年生，$B_2=$2年生，$B_3=$3年生とします。

ここでランダムに一人を選び出した時，それが「3年男子」である確率はどれぐらいでしょうか？ これは「学年は3」かつ「性別は男」なので，学年と性別の両方の条件が入った同時確率といい，$P(A_i,B_j)$と書きます。これは表から，次のように計算できますね。
\[ P(A_3,B_1)= \frac{13}{100}=0.13 \]

次に用語として\keyterm{周辺確率}{しゅうへんかくりつ}{marginal probability}というのを考えます。これはある分割に関して全て足し上げたもの，というものですが，要するに表\ref{tbl::24_02}の合計欄のことです。例えば「男子生徒についての周辺確率」は，男子生徒が学年によって3つに分割されていますが，これを全て足し上げるので，1年男子+2年男子+3年男子で求まるものです。
式で言うと，次のような意味になります。

\[ \sum_{i=1}^3 p(A_i,B_j) = p(B_j) \]
\[ \sum_{j=1}^2 p(A_i,B_j) = p(A_i) \]

全部の数字を足し合わせるので，その変数の確率の話ではなくなる，と思ってください。このように一方の変数の全ての確立を足し合わせることを\textbf{周辺化}する，といいます。

では次のステップ。今度は\keyterm{条件つき確率}{じょうけんつきかくりつ}{conditional probability}というのを導入します。無作為に学生を選び出す時，「選ばれたのが女子である」ということが先にわかったとします\footnote{例えば名前とか容易に類推できた，といったようなシーンです。}。その時，その生徒が2年生である確率はどれくらいでしょうか？ この問いを「女子という条件がついた上での確率」という意味で，条件つき確率といい，表現としては$P(\mbox{2年生}|\mbox{女子})$と書きます。2つの分割を縦棒で区切り，条件の方が後ろに書きます。「女子という条件が与えられた上での二年生の確率は」というように読んだりします。

ちょっと状況はややこしくなりましたが，確率の計算そのものはそれほど難しくありませんね。女子である，というのが決まっているので，分母が変わるだけなのです。つまり，
\[P(A_2|B_2)=\frac{20}{60}=0.333... \]
となります。

条件つき確率の計算は分母が変わったと思えばいいのですね。この条件つき確率は，同時確率を周辺確率で割っても計算できます。すなわち，
\[ P(A_2|B_2) = \frac{P(A_2,B_2)}{P(B_2)}=\frac{0.2}{0.6}=0.333... \]
となるわけです。

今回は性別という条件が与えられた時に，学年はどうなるかという条件をつけた確率を考えました。性別が一年生の時，二年生の時，三年生の時，それぞれで分母が変わりますし，条件付き確率も変化します。このように，性別が変わると学年の確率が変わる，という関係にあるので，この性別と学年という二つの変数は関係しあっていることがわかります。逆に，一方の変数が変わっても他方の変数に影響しない場合，両者は\keyterm{独立}{どくりつ}{independent}だ，といます。
確率において独立であるとは，全ての$i,j,k$において次の関係が成り立つことを言います。

\[ p(A_i| B_j) = p(A_i|B_k) \]

これまでも(確率的に)独立である例の話はいろいろしてきました。例えばコイントスを続けて行うときは，それぞれの試行は独立だよね，という話をしました(Pp.\pageref{iidexample})。一回のコイントスの結果が次のコイントスの結果に影響しない，という意味です\footnote{逆に一回のコイントスが次のコイントスに影響するとはどういうシーンでしょうか? 無理矢理考えるなら，コインが意思を持って「さっき表になったから，次は裏にしようとする」とか，トスしたコインを受け取り損ねて地面に落下し，片方の面が欠けて他方の面が出やすくなった，というようなシーンです。コインは意思を持ってないし，硬いものですから，これらはやはり考えにくい話ですよね。}。

ここで独立しているのならば，ここまでの式展開から，

\begin{equation}
    \begin{array}{ccc}
        p(A_i| B_j) &=& p(A_i|B_k) \\
        \frac{p(A_i,B_j)}{p(B_j)} &=& \frac{p(A_i,B_k)}{p(B_k)}\\
        p(A_i,B_j)p(B_k) &=& p(A_i,B_k)p(B_j) \\
        &&\mbox{すべての$k$の状態について$B_k$を数え上げる，つまりBを周辺化します}\\
        p(A_i,B_j)\sum_k=1^K p(B_k) &=& p(B_j)\sum_{k=1}^K p(A_i,B_k) \\
        p(A_i,B_j) &=& p(B_j)p(A_i)
    \end{array}
\end{equation}

ここからわかるように，AとBが独立である場合は，それぞれの確率の積で表現できます。これまでも，コイントスやサイコロの例で，2回続けて計算する場合はその積で表現してきましたよね。
確率$\theta$で表が出るコインで，2回続けて表が出る場合は$\theta \times \theta=\theta^2$，表・裏と出る場合は$\theta \times (1-\theta)$といった表現をしてきました。
独立した事象は掛け算で表せる，というのはこういう理屈からなんですね。


さて，これで用語の準備は終了です。ここから式の展開が面白くなってきます。

あらためて，条件付き確率は同時確率を周辺確率で割ったもの，でしたね。つまり，
\[ P(B_j | A_i) = \frac{P(B_j,A_i)}{P(A_i)} \]
ここで右辺の分母を両辺にかけて，
\[ P(A_i)P(B_j|A_i) = P(B_j,A_i) \]
とします。これはA,Bが何であるかにかかわらず，一般的に成立します。たとえば$P(A_i|B_j)$と記号をひっくり返しても，
\[ P(A_i|B_j) = \frac{P(A_i,B_j)}{P(B_j)} \]
より
\[P(B_j)P(A_i|B_j) = P(A_i,B_j) \]
としても構いませんね。


ところで，「2年の男子」と「男子の2年」というのは同じ意味です。つまり，
\[P(B_j,A_i) = P(A_i, B_j) \]
ですから，上の式をつなぎ合わせて
\[P(A_i)P(B_j|A_i) = P(B_j)P(A_i|B_j) \]
とできます。ここで両辺を$P(B_j)$で割って清書すると，
\[ P(A_i | B_j )=\frac{P(B_j|A_i)P(A_i)}{P(B_j)} \]
が得られます。やったね！ベイズの定理が導出できました。


さて，これの何がすごいのでしょう？ 数式を「記号の羅列だ！」と思って拒否するのではなく，その意味をしっかり味わってみましょう。

\section{ベイズの定理の意味}
\subsection{因果の向きについて}
さてベイズの定理のすごいところは，\underline{左辺と右辺の中にある条件つき確率の中身がひっくり返っているところ}にあります。
これがひっくり返ると何がうれしいのでしょうか。
ベイズの定理の左辺には$P(A|B)$があり，右辺は分子に$P(B|A)$を含んでいます。$P(A)$(分子の2つ目の項目)や$P(B)$(分母)のことは一旦忘れてくれてもいいですし，
\[ P(A|B) = P(B|A) \times \frac{P(A)}{P(B)}\]
と考えて，$P(A|B)$が$P(B|A)$の関数である，と読み取ってくれても構いません。

ここでAを「心のモデル」，Bを「データ」と考えてみてください。
私たちは心の仕組みを考える時に，モデルをたてます。「あの人が$\diamond\diamond\diamond\diamond$な行動をとったのは，心の状態が$\dagger\dagger\dagger\dagger$だったからではないか」というのは心というモデルを立てて考えるということなのでした。心の状態が$\dagger\dagger\dagger\dagger$であれば，人は$\diamond\diamond\diamond\diamond$な行動をとるであろう，というのが研究者の仮説，モデルなのです。つまり，$P(\mbox{データ}|\mbox{モデル})$が心理学でやっていることなのです。

ところが実際には，手に入れることができるのはデータだけです。心がどうなっているかはわかりませんが，人が$\diamond\diamond\diamond\diamond$な行動をとった，という事実だけがそこにあるのです。この事実を使って，我々のモデルがどれぐらい正しいのか。あるいはモデルのパラメータはどうなっているのか，を検証したいわけです。つまり，研究の営みとは$P(\mbox{モデル}|\mbox{データ})$であり，データを所与(条件)としたときの，モデルの確率を考えることなのです。データをとってモデルを検証するという営みが成り立つ根拠を，ベイズの定理が示してくれていると言えるかもしれません。大事なのは，不確かな状況においては我々は確率的にしか判断をできず，ベイズの定理では条件の向きをひっくり返して計算できる，という点です。原因が結果を引き起こす確率が知りたい時に，結果を手にすることで，原因がどれほどあり得たかを考えることができる。それがベイズの定理のポイントなのです。ベイズの定理は因果の向きを逆にすることができる，ということを示しているので，\keytermL{逆確率の法則}{ぎゃくかくりつのほうそく}と呼ばれることもあります\footnote{ベイズ自身がこの法則を発見した時に，「なんかおかしなことになってるかも」と思って結果を公表しなかったと言われています。神様が決めた世界を裏返しにすることに，背徳感のようなものを感じたのかもしれません。}。

\subsection{確率という数字，再考}
ここまでベイズの定理の説明は，性別や学年など離散的な数字をもとに説明してきましたが，連続的な確率変数でも成り立ちます。連続的な確率変数の場合，全ての確率を足し合わせるという操作が積分することに変わりますが，それ以外は本質的に変わりません。一般的に書くと，

\[ f(\bm{\theta}|\bm{x}) = \frac{f(\bm{x}|\bm{\theta})f(\bm{\theta})}{f(\bm{x})} \]

となります。ここで$\bm{\theta}$はパラメータで，複数ある場合があるのでベクトルの表記(太字)にしてあります。$\bm{x}$はデータのことで，これも一つのデータではないので太字にしています。$f$は連続的な確率関数だと思ってください。
先ほど説明したように，この関数は

\[ f(\bm{\theta}|\bm{x}) = f(\bm{x}|\bm{\theta}) \times \frac{f(\bm{\theta})}{f(\bm{x})} \]

と考えても構いません。ここで$f(\bm{x}|\bm{\theta})$は，パラメータが所与の元でのデータの出てくる確率ですから，\keytermL{尤度}{ゆうど}だ，ということがわかります。式の左辺は$ f(\bm{\theta}|\bm{x})$ですから，データが所与の時のパラメータの確率ですね。これを特に\keyterm{事後確率}{じごかくりつ}{posterior probability}と言い，それに対して左辺の$f(\bm{\theta})$は\keyterm{事前確率}{じぜんかくりつ}{prior probability}といいます。
ベイズの式は，確率を確率に変換する公式なのです。

ところでこの\keytermL{確率}{かくりつ}って何なのでしょうね。パラメータについての事前の確率とは，一体どういう意味があるのでしょうか。

思い出してほしいのですが，第\label{chapter7} 章で確率の公理を導入しました。その時説明したように，確率とは次の条件を満たす数字だということでした。
\begin{enumerate}
    \item 確率の値は非負（正または0の数字）
    \item 全標本空間における全ての事象の確率の合計は1.0
    \item 任意の2つの相互に排他的な事象について，どちらか一方が発生する確率は個々の確率の和である
\end{enumerate}
そして，確率には\keytermL{頻度主義的}{頻度主義的}な考え方と，\keytermL{主観確率}{主観確率}という考え方があって，どちらでもこの公理を満たしている限り確率として計算できます，という話をしてきました。

高校数学までは頻度主義的な意味で確率が説明されてきたと思います。すなわち，コインやサイコロを\textbf{無限に}投げ続けると，パラメータの値は1/2や1/6に近づいていくだろう，ということです。この無限に繰り返すと一つの値に収束するというのは，確率的収束という数学的な証明がされる考え方ですが，現実には無限に行うというのはあり得ません。心理学の実験も，同じ実験を無限に繰り返すと，という設定のもとで考えて帰無仮説検定を行ったりしますが，それがそもそもあり得ない前提じゃないか，と言われたらどうしようもないのです\footnote{さらに帰無仮説というあり得ない仮定を置いてますから，一体どこの世界の話をしているのやら，という気持ちになります。}。無限の回数を考えて，その相対頻度，割合で確率を考えるので頻度主義と呼ばれています。

これに対して天気予報などは，新年の確からしさを確率で表現しています。明日晴れる確率が30\%というのは，明日という日が無限に繰り返されればそのうちの3割は晴れている，と言うのはやはりおかしいわけです。どんなループに巻き込まれてしまったのでしょうか，明日が無限に繰り返されるなんて！
そういう意味では，普段私たちが慣れ親しんでいるのは，自分達の確信の強さ，どの程度その事象を信じているかと言う信念の強さを表している確率なのです。明日100\%雨が降る，というのは完全にその状態を信じているということですし，50\%の降水確率は雨が降る可能性が半分，降らない可能性が半分だと考えていると言うことですね。

相対頻度・割合で考えるときは，全体で割った平均的な数字を出すことになります。この頻度主義的な考え方ですと，例えば野球のシーズン最初の試合，勝ったチームは勝率100\%になります。一試合しかしていないので，1/1だからです。でも誰も「このチームは今シーズン一度も負けないだろう」と本気で信じているわけではないですよね。あるいは最初の打席でヒットが出なかったからといって，打率0割，永遠に打てない打者だとは誰も考えません。本当はもう少し低い勝率，高い打率だと思っているのに，そう表現できていないのです。

% \footnote{注意が必要なのは，モデルの正しい確率が検証されるわけではないところです。モデルを仮定すると，そのモデルから考えられる結果の確率がこれぐらいになる，という推論をしているのであり，モデルがそもそも前提されていますから，モデルが間違っていれば(例えばコイントスという現象に正規分布を当てはめて考えるとか)，その推論は全然正しいものではありません。}。





ベイズの定理を使ってパラメータを推論する時に，この定理が根本にあるということの意味が伝わりましたでしょうか。



% \item[ベイズの定理] ベイズの定理について，数式を様々な角度から眺めてその意味を理解する。同時確率，周辺確率の定義などを用いて，条件付き確率の式を変形していくことができ，最終的に右辺と左辺で条件が逆になる場合の確率が計算できることがわかる。これがベイズの定理と呼ばれるものであり，逆確率の定理，法則と呼ばれることもある。ベイズの定理は，事前確率に尤度と周辺尤度の比をかけたものとして理解することもできるし，尤度関数を確率として捉えるために変形させる項として事前確率と周辺尤度の比をかけたもの，と考えることもできる。
% \begin{flushright}
%     $\to$ \citet{Maeda2017}のP.93--109.尤度を確率分布に変えることの分かりやすい説明としては，\citet{Lambert201805}のTable4.1.Xが良い。
% \end{flushright}

% \item[ベイズの定理を使った推論の例]ベイズの定理について，意味的な解釈を膨らませておこう。ベイズの定理は確率のアップデートとして考えることができる。我々が推論するというのは，確率をアップデートすることだと考えることもできる，ということについてはすでに触れており，確率という数字が従うべきルールさえ守られていれば，確率は人間の自然な推論モデルとして考えることができるだろう。ただし頻度主義的な確立の理解と，主観確率による理解の違いには改めて注意が必要である。またベイズの定理を使っての推論の例として，病気の罹患率の例とモンティ・ホール問題を取り上げる。これらのポイントは，「事前分布がわかっているということ」と「情報がアップデートされる」ということを意識するかどうかにある。
% \begin{flushright}
%     $\to$ \citet{Maeda2017}のP.15--23.あるいは\citet{豊田201506}のP.8--12.
% \end{flushright}

\ssection{ベイズの定理を使った推論の例}

コロナウイルスに感染しているかどうかは，PCR検査で明らかになりますが，これらの検査の結果というのも確率的な要素があるのはご存知でしょうか。本当は感染していないのに，検査では陽性になってしまう，ということがゼロではありません。こういうのを偽陽性といいます。逆に感染しているのに，検査で検出できずに陰性になる，偽陰性ということもあります。いずれも確率的な話です。

コロナに感染している$\to$検査で陽性になる，というのが因果関係ですが，我々は実際には検査で陽性である$\to$コロナに感染しているかもしれない，というように推論します。これを今日学んだ確率の言葉で書いてみましょう。

感染している状態で，陽性反応が出るというのは条件つき確率として表現できますね。$P(\mbox{陽性}|\mbox{感染})$，というわけです。さてここに，そもそもの感染率というのが別途見積もれたとします。$P(\mbox{感染})$ですね。さらにもうひとつ。感染しているかどうかに関係なく，陽性であるという結果だけ取り上げて$P(\mbox{陽性})$を考えることもできます。これが揃うと，先ほどの定理を使って
\[ \frac{P(\mbox{陽性}|\mbox{感染})P(\mbox{感染})}{P(\mbox{陽性})} = P(\mbox{感染}|\mbox{陽性})\]
を計算できます。つまり陽性になったという結果から，本当にかかっているかどうかの確率を計算できるのです。

多くの場合，検査で陽性だったら「もうだめだー」と思ってしまいがちです。これは陽性であると確実に検出できる，偽陽性がゼロの検査であれば正しい判断です。が，実際には確率的に判断すべき問題です。具体的に数字を例にして考えてみましょう。

\begin{itemize}
    \item 40歳の女性が乳がんにかかる確率は1\%です。
    \item 乳がんである人が乳房X線検査で陽性と出る確率は90\%です。
    \item 本当は乳がんでなくても，検査で陽性と出る確率は8\%にすぎません。
    \item あなたが40歳の女性で検査結果が陽性だったとしたら，あなたが乳がんである確率はどれぐらいでしょうか。
\end{itemize}
これは先ほどの公式に当てはめて考えると，
\[\frac{0.9\times0.01}{0.9\times0.01+0.08\times0.99}=0.102 \]

となり，10\%ちょっとだということになります\footnote{おっと，計算についていけないぞ！と思った人。大丈夫です，いい方法があります。確率の計算は小数点を含みますし相対的な度数なのでわかりにくいですから，絶対的な度数に戻して考えればいいのです。例えば10000人の世界だとして考えると，乳がんの人は1\%=100人です。乳がんでない人は9900人です。陽性になる人は乳がんで陽性反応が出るひと($100\times0.9=90$)と乳がんじゃないのに陽性反応が出る人($9900\times0.08=792$)の和で，882人です。陽性反応が出た人の中で本当に乳がんであるひとは，$90/882=0.102$となります。}。ずいぶんと印象が変わりますよね。PCR検査を複数回行うのは，これと同じような背景があるからです。

\section{課題}
\end{document}
